%%%%%%%%%%%%%%%%%%%%%%%%%%%%%%%%%%%%%%%%%%%%%%%%%%%%%%%
%
%           Thesis Main File
%
%                                   author:Shouta Samejima
%                                   date:19/01/13
%
%                                   edit:Koki Nakamura
%                                   date:23/12/21
%
%                                   edit: Masaki Tanaka
%                                   date:24/02/09
%%%%%%%%%%%%%%%%%%%%%%%%%%%%%%%%%%%%%%%%%%%%%%%%%%%%%%%

\documentclass[a4j,11pt,onecolumn,twoside,titlepage,openright,final, openany]{jreport}
%\documentclass[a4j,11pt,onecolumn,twoside,titlepage,openright,final]{jreport}

%***************************************** ＜Preamble＞ ******************************************%
\usepackage{graphicx}
\usepackage[dvipdfmx]{color}

\usepackage{algorithm} %algorithmicやalgorithmicxなどの親パッケージ
\usepackage{algpseudocode} %疑似コード作成：Functionの使用

\usepackage{amsmath, amssymb} %数式表現の拡張
\usepackage{type1cm} %実数Rとかをかっこよく

\usepackage{bm} %太文字の使用
%\usepackage{booktabs} %表中の罫線の種類の追加
%\usepackage{cases} %場合分けの使用
\usepackage[strict]{changepage} %奇数・偶数ページのレイアウト変更
\usepackage{comment} %複数行コメントアウトの仕様
\usepackage{empheq} %amsmathの拡張
\usepackage{epsfig} %eps形式の図の使用
\usepackage{fancyhdr} %ヘッダ・フッターのカスタマイズ：fancyスタイルの使用可になる
\usepackage[T1]{fontenc} %フォントエンコーディングをT1に変更
\usepackage{here} %図，表を強制的に配置
\usepackage{latexsym} %数式で使用可能な記号の拡張
\usepackage{listings} %ソースコードの記述
%\usepackage{longtable} %表機能の拡張：ページを跨いだ表の作成
%\usepackage{makeidx} %索引の作成
%\usepackage{multirow} %縦方向のセルの結合
\usepackage{oldgerm} %古いドイツ語フォントの使用
\usepackage{pgfplots} %高度なグラフの描画
\usepackage{physics} %物理系の数式の記述
%\usepackage{slashbox} %表で斜線を含んだラベルセルの使用
\usepackage{subfiles} %プリアンブルを共有したファイル分割
%\usepackage{tabularx} %表組の幅指定や書式指定
\usepackage{textcomp} %テキストシンボルの使用
%\usepackage{tikz} %関数のグラフの記述
\usepackage{verbatim} %入力通りにテキストを出力
\usepackage{amsmath}


%----------------------- <Pass setting> -----------------------%
\graphicspath{{fig/}} %{}内にmain.texから画像が入っているフォルダまでのpathを記述

\providecommand{\subinclude}[1]{\subfileinclude{chapter/#1}} %サブtexが入っているフォルダまでのpathを記述
\providecommand{\supinput}[1]{\input{supplement/#1}} %補助texが入っているフォルダまでのpathを記述
\providecommand{\supinclude}[1]{\include{supplement/#1}} %補助texが入っているフォルダまでのpathを記述
\providecommand{\subinclude}[2]{\subfileinclude{chapter/#2}} %サブtexが入っているフォルダまでのpathを記述
\providecommand{\supinput}[2]{\input{supplement/#2}} %補助texが入っているフォルダまでのpathを記述
\providecommand{\supinclude}[2]{\include{supplement/#2}} %補助texが入っているフォルダまでのpathを記述
\providecommand{\subinclude}[3]{\subfileinclude{chapter/#3}} %サブtexが入っているフォルダまでのpathを記述
\providecommand{\supinput}[3]{\input{supplement/#3}} %補助texが入っているフォルダまでのpathを記述
\providecommand{\supinclude}[3]{\include{supplement/#3}} %補助texが入っているフォルダまでのpathを記述
\providecommand{\subinclude}[4]{\subfileinclude{chapter/#4}} %サブtexが入っているフォルダまでのpathを記述
\providecommand{\supinput}[4]{\input{supplement/#4}} %補助texが入っているフォルダまでのpathを記述
\providecommand{\supinclude}[4]{\include{supplement/#4}} %補助texが入っているフォルダまでのpathを記述
\providecommand{\subinclude}[5]{\subfileinclude{chapter/#5}} %サブtexが入っているフォルダまでのpathを記述
\providecommand{\supinput}[5]{\input{supplement/#5}} %補助texが入っているフォルダまでのpathを記述
\providecommand{\supinclude}[5]{\include{supplement/#5}} %補助texが入っているフォルダまでのpathを記述



%------------------------------------------------------------%



\supinput{preamble}

%***************************************** ＜Main body＞ ********************************************%
\begin{document}

\supinput{title} %表紙



%***** <Customizing header and footer> *****%
\pagestyle{fancy}
\lhead[]{} %[偶数ページ]{奇数ページ}
\chead[目 次]{目 次} %\leftmark=デフォルトは節見出し：\rightmark=デフォルトは章見出し
\rhead[]{}

\cfoot{\thepage} %[ ]を省略すると{ }で指定された内容が偶数ページ・奇数ページの両ページに適用される

\renewcommand{\headrulewidth}{0.4pt} %ヘッダ線(上線)の長さ 0.4ptがデフォルト
\renewcommand{\footrulewidth}{0pt} %フッタ線(下線)の長さ 0ptがデフォルト

\makeatletter
\def\chaptermark#1{\markboth {\ifnum \c@secnumdepth>\m@ne\@chapapp\ \thechapter \@chappos\hskip 1zw \fi #1}{}} %chapterコマンドが\markbothに渡す引数のカスタマイズ
\def\sectionmark#1{\markright {\ifnum \c@secnumdepth >\z@ \thesection \hskip 1zw \fi #1}}%\sectionコマンドが\markrightに渡す引数のカスタマイズ
\makeatother
%***************************************%



%***** <Setting page numter> *****%
\pagenumbering{roman} %目次用のページ番号付け
\setcounter{page}{1} %ページ番号をゼロに設定
%******************************%



%***** <Making contents> *****%
\tableofcontents %目次の作成
\listoffigures %図目次
%\chead[図 目 次]{図 目 次}
%\listoftables %表目次
%\chead[表 目 次]{表 目 次}
%***************************%



%ヘッダの再定義
\chead[\leftmark]{\rightmark}

%***** <Reading each chapters> *******************************************************************%
% \subinclude{chapter1/chapter1.tex}
\subinclude{chapter2/chapter2.tex}
% \subinclude{chapter3/chapter3.tex}
% \subinclude{chapter4/chapter4.tex}

% \supinclude{thanks}
%\cleardoublepage %ページ合わせの空白ページ挿入

%\chead[参考文献]{参考文献}
\supinclude{reference}
% \supinclude{publications}

%\chead[付録]{付録}
%\include{Appendix}
%\cleardoublepage
%\include{summary}
%\cleardoublepage

%\bibliographystyle{jplain} %bibtex使用ならコメントアウト解除
%\bibliography{}    %\bibliography{bibファイル名}

\end{document}